\begin{abstract}
Electrocardiogram (ECG) foundation models are deployed for cardiac risk stratification, yet
existing benchmarks measure discrimination exclusively via AUROC, leaving probability
calibration invisible. A model that ranks patients perfectly but assigns systematically inflated
or deflated probabilities can cause direct clinical harm when its outputs are used to drive
treatment thresholds. We present the first systematic calibration audit of ECG foundation
models: six encoders---five ECG foundation models (MERL, ST-MEM, HuBERT-ECG, ECGFM-KED,
ECG-FM) and a supervised S4D baseline---across four datasets, totalling 24 model$\times$dataset
cells (23 evaluated) over 18\,877 ECG records (2\,000 PTB-XL, 6\,877 CPSC2018, 8\,000
CODE-15\%, 2\,000 MIMIC-IV-ECG) under patient-level leakage-controlled splits, with bootstrap
confidence intervals on every headline metric.

Our central finding is a \textbf{discrimination--calibration tradeoff} on the multi-label
PTB-XL rhythm task: ST-MEM achieves the highest AUROC (0.838) but poor calibration
(ECE 0.124), HuBERT-ECG is the worst calibrated (ECE 0.136), while the supervised S4D
baseline is best calibrated (ECE 0.064) despite the lowest AUROC (0.771). All five
foundation models exceed S4D's ECE by 0.015--0.072. Strikingly, this tradeoff
\emph{disappears} on the lower-co-occurrence endpoint tasks---CPSC2018 (ECE 0.023--0.051),
CODE-15\% (ECE 0.010--0.032), and (report-labelled) MIMIC-IV-ECG (ECE 0.043--0.106)---where
high-AUROC models are also well-calibrated. On MIMIC-IV-ECG, MERL attains the best discrimination
(AUROC 0.861); critically, ECG-FM, which was \emph{pretrained on MIMIC-IV-ECG}, reaches
only 0.771---behind MERL and ST-MEM (in all 10 resplits) and no better than the supervised
S4D baseline (0.793)---demonstrating that within-domain pretraining does not guarantee
superior linear-probe performance. We link
the PTB-XL tradeoff to task structure: multi-label correlated classes create geometrically
harder representation-to-probability mappings.

Post-hoc isotonic recalibration reduces ECE by 43--71\% on PTB-XL (ST-MEM:
$0.124\!\to\!0.044$; HuBERT-ECG: $0.136\!\to\!0.040$) while causing only modest AUROC
reductions of 0.5--1.7\%. A graded pretraining-objective signature emerges: MERL
(contrastive InfoNCE) shows the strongest underconfidence ($\text{SCE}=-0.003$),
masked-prediction models (ST-MEM, HuBERT-ECG) are near-neutral, and knowledge-enhanced,
hybrid, and supervised models are mildly overconfident ($\text{SCE}=+0.007$--$+0.008$),
consistent with documented underconfidence of contrastive foundation models. Cross-dataset
analysis reveals that the same model's ECE varies by up to roughly 6-fold across datasets,
underscoring that calibration must be audited per deployment context.

We introduce \emph{PhysioCalib}, a severity-weighted isotonic calibration method, and
open-source \texttt{ecg-audit}, a pip-installable, unit-tested Python package for
one-command calibration auditing, recalibration, and decision-curve analysis. Our findings
demonstrate that probability calibration should be a first-class evaluation metric alongside
AUROC in ECG AI benchmarks.
\end{abstract}
